\chapter{B\_trajectory\_conversion.py}
\label{ch:scriptB}

This script converts the raw trajectory file obtained from the total station measurements into the format which is compatible by Pulsalyzer software. This converted trajectory is also used as the trajectory input for the Python workflow explained in \autoref{ch:pythonWorkflow}.

\section{File Structure}
The measurements captured by the total station is logged by the Allterra TS Logger software in a .txt file format. Each row in this file contains a time-of-day timestamp (HHMMSS.SSS format in UTC), previously set station coordinates, a horizontal angle (azimuth), a vertical angle (zenith), and a slope distance.

This script reads the raw trajectory file and performs the following operations:

\begin{enumerate}
    \item Recomputes Cartesian (X, Y, Z) coordinates of tracked points from the polar coordinates (azimuth, zenith, slope distance) with respect to the total station as origin.
    \item Converts from left-handed to right-handed coordinate system (negates Y).
    \item Applies a lever arm correction along Z to account for the lever arm offset between the tracked prism and the ULi system.
    \item Converts time-of-day timestamps to seconds since UNIX epoch (1 January 1970 00:00:00 UTC) and adjusts for the UTC-to-TAI time standard difference.
    \item Outputs the trajectory data in tab-delimited Pulsalyzer format.
\end{enumerate}

\section{User-Configurable Parameters}

\begin{table}[ht!]
    \centering
    \begin{tabularx}{\textwidth}{llX}
        \toprule
        Variable                       & Default            & Description                                                                               \\
        \midrule
        Year, Month, Day               & '2025', '08', '06' & Date of data collection                                                                   \\
        x\_rot, y\_rot, z\_rot         & 0, 0, 0            & Scanner orientation (in degrees)                                                          \\
        utc2tai\_offset                & 37                 & UTC-to-TAI offset in seconds                                                              \\
        prism\_to\_scanner\_lever\_arm & 1.65               & Vertical distance (m) from prism to ULi lens                                              \\
        input\_trajectory\_file        & R"..."             & Full path to raw .txt trajectory file (e.g.: \texttt{\detokenize{R"C:\Folder\file.txt"}}) \\
        \bottomrule
    \end{tabularx}
\end{table}

\newpage
\section{Output File}
The output is a tab-delimited \texttt{.txt} file containing the processed trajectory with the following columns and a corresponding header:

\begin{table}[ht!]
    \centering
    \begin{tabularx}{\textwidth}{llX}
        \toprule
        Column         & Unit                     & Description                                                                              \\ \midrule
        \texttt{time}  & seconds (s)              & Timestamps since UNIX epoch (including UTC-to-TAI offset)                                \\
        \texttt{x}     & metres (m)               & Cartesian X coordinate in the right-handed total station frame                           \\
        \texttt{y}     & metres (m)               & Cartesian Y coordinate in the right-handed total station frame                           \\
        \texttt{z}     & metres (m)               & Cartesian Z coordinate in the right-handed total station frame with lever arm correction \\
        \texttt{roll}  & degrees (\unit{\degree}) & Scanner roll orientation (Default \qty{0.00}{\degree})                                   \\
        \texttt{pitch} & degrees (\unit{\degree}) & Scanner pitch orientation (Default \qty{0.00}{\degree})                                  \\
        \texttt{yaw}   & degrees (\unit{\degree}) & Scanner yaw orientation (Default \qty{0.00}{\degree})                                    \\
        \bottomrule
    \end{tabularx}
\end{table}

The output file is saved in the same directory as the input trajectory file with the following naming convention:

\begin{verbatim}
converted_{Year}_{Month}_{Day}_offset_{utc2tai_offset}_traj_{x_rot}_{y_rot}_{z_rot}.txt
\end{verbatim}

where,

\begin{itemize}
    \item \texttt{Year, Month, Day} are the date of data collection
    \item \texttt{utc2tai\_offset} is the UTC-to-TAI offset in seconds
    \item \texttt{x\_rot, y\_rot, z\_rot} are the scanner orientation offsets
\end{itemize}

\section{Example of conversion}

\subsection*{Raw trajectory}

\begin{lstlisting}[numbers=none]
101533.231 548775.866 5937527.313 10.950 208.58550 99.53554 35.764
101533.281 548775.866 5937527.312 10.950 208.58530 99.53544 35.763
101533.331 548775.866 5937527.313 10.950 208.58550 99.53534 35.764
101533.382 548775.867 5937527.313 10.950 208.58570 99.53534 35.764
101533.433 548775.866 5937527.313 10.950 208.58560 99.53534 35.764
\end{lstlisting}

\subsection*{Converted trajectory}

\begin{lstlisting}[numbers=none]
time	x	y	z	roll	pitch	yaw
1754475370.231000	-35.438322	4.808424	-1.389078	0.000	0.000	0.000
1754475370.281000	-35.437346	4.808178	-1.389029	0.000	0.000	0.000
1754475370.331000	-35.438321	4.808424	-1.388966	0.000	0.000	0.000
1754475370.382000	-35.438306	4.808535	-1.388966	0.000	0.000	0.000
1754475370.433000	-35.438314	4.808479	-1.388966	0.000	0.000	0.000
\end{lstlisting}